\section{The Gate, and What It Needs}
\label{sec:ch16-the-gate-and-what-it-needs}

\index{schema check!the tool that exists|(}%
Section~\ref{sec:ch16-what-the-composer-is-checking} ends on a question the
composer cannot be asked, so ask the tool that was built to answer it. Cosmo
ships one, \code{wgc subgraph check}, a subcommand of the same binary this
book has composed with since chapter~\ref{ch:composition}. Here is what it
does on this machine:

\begin{minted}[fontsize=\footnotesize,breakanywhere]{text}
Organization slug is not set. Please run `wgc auth login` to set the organization slug.
\end{minted}

\index{schema check!needs a control plane}%
Exit one. The same message comes back from \code{wgc proposal create}, from
\code{wgc subgraph fix} and from
\code{wgc federated-graph check}.
That is not a missing flag. The command's own help says what it is
doing, and the giveaway is the last four words: it checks for breaking changes
and composition errors \emph{across all connected federated graphs}. Knowing
which graphs are connected, and what schema each of them is currently serving,
is what a control plane is for, and \code{wgc router compose} has been the
exception rather than the rule all along. Its help says so too, in a sentence
this book has been quietly relying on since
chapter~\ref{ch:composition}: it composes locally
\enquote{without a control-plane connection}.

\subsection{What the Check Is Actually Comparing}

\index{schema check!two comparisons, not one}%
The interesting half is not that it needs a server. It is what the server is
holding. A schema check makes two comparisons. The first is a diff, between
the candidate document you hand it and the one the control plane already has
for that subgraph, across every graph that subgraph is part of, which is what
the help text means by composition errors in the plural. The second is the one
nothing local can do. It takes the changes the first found and asks whether
any client actually used them.

\index{schema check!the traffic comes from the router}%
Where the traffic comes from is the router. Cosmo's documentation is plain
about it and about what follows~\autocite{cosmochecks}:

\begin{quote}
This is done by sending schema usage traffic to Cosmo Cloud from your routers.
If you propose a breaking change and no active clients use the affected schema
changes the check will pass.
\end{quote}

\index{schema check!seven days by default}%
The window is seven days by default and is a namespace setting. So a field
nobody has selected in a week is a field you may delete, and a client that
runs a report on the first of the month is a client the check has never heard
of. That is a policy dial rather than a defect, and it is one worth setting
deliberately, because the default is a week and quarterly things exist.

\index{schema check!directives are not tracked}%
One limit in the same document deserves more attention than its two sentences
suggest, given how much of this book's graph is directives:

\begin{quote}
Currently, we do not track the usage of directives. This means that changing or
deleting directives will always be detected as a non-breaking change, even if
clients are using them.
\end{quote}

\index{schema check!the option that turns the traffic off}%
And there is a flag for the case where you do not want the second comparison
at all, which is the honest way to read what the first one is worth on its
own: \code{--skip-traffic-check}, described in the help as skipping the check
for client traffic, after which any breaking change fails the run. That is
the setting for a graph whose clients you cannot enumerate, and it converts
the tool from a question about your users into a question about your schema.

\index{schema check!Apollo's is the same shape}%
Apollo's is the same tool with different nouns. \code{rover subgraph check}
needs an API key, a schema already published to their service, and a graph
ref, which is their name for the identifier of one graph and one variant of
it. It diffs against the published schema and then against operations executed
inside a window that is also seven days by default. Its categories are keyed to usage rather than to
shape, so its own definition of a removed field is a field \emph{used by at
least one operation}. And its reference page concedes something I have not
seen a vendor concede often~\autocite{apollochecksref}:

\begin{quote}
In some cases, in-place updates are compatible with affected clients at runtime
(such as a type rename or a conversion from an object to an interface that uses
the same fields). However, schema checks still marks these as breaking changes,
because validation does not have enough information to ensure that they are
safe.
\end{quote}

\index{schema check!false positives on purpose}%
A check that reports safe changes as breaking, on purpose, because it cannot
tell. I would rather have that than the alternative, and it is worth knowing
before the first argument about a check that is wrong.

\subsection{Why This Book Does Not Run One}

\index{control plane!what self-hosting it costs}%
So I should say why this book has no such report in it, rather than leave the
absence to be read as an oversight. Cosmo is open source and the control plane
can be self-hosted. What self-hosting it means, in WunderGraph's own
documentation, is Docker Compose on a laptop or Kubernetes with a Helm chart
in production, standing up a Node service, PostgreSQL, Redis, ClickHouse, an
S3-compatible object store, Keycloak with a PostgreSQL of its own, a CDN
server, an OpenTelemetry collector and a metrics
collector~\autocite{cosmoselfhosted}. This book's running example has four
processes, four SQLite files and one downloaded binary, and it has no
container in it anywhere, which was settled before
chapter~\ref{ch:do-you-need-this} was written and has held ever since.

\index{control plane!the honest position}%
So I am not going to stand one up to write a chapter, and I am not going to
describe a report I have not seen. What I can tell you is what the command
needs, which is measured, and what it compares, which is documented. What this
chapter does not contain is Cosmo's own verdict on
chapter~\ref{ch:null-blast-radius}'s five fields, because nobody ran it.

\subsection{The Half You Can Build Without One}

\index{schema check!compose twice and diff}%
The first comparison needs no server at all, and
section~\ref{sec:ch16-the-change-this-book-already-shipped} is it, performed
as a procedure rather than as a demonstration. Compose the four documents you
have published. Compose the four you are about to publish. Pull
\code{engineConfig.graphqlSchema} out of each config and diff them. Two
invocations of a command you already run and one diff, and the report is the
five changed lines that section prints.

\index{schema check!diff the composed schema, not the subgraph}%
Diff the composed schema rather than the subgraph document, and the reason is
in those five lines. Six moved across the four documents and five fields
moved in the supergraph, because \code{Query.nodes} is declared twice and
merged once. A per-subgraph diff would have reported six changes to a schema
that has five, and it would have missed the general case entirely: a field can
change in the composed schema because a \emph{different} subgraph changed. The
supergraph is what a client reads, so the supergraph is what you diff.

\index{schema check!why the vendor takes one document and you compose two}%
That is not a disagreement with \code{wgc subgraph check}, which takes one
subgraph document and nothing else. It can, because a control plane is holding
the other three and the schema they last composed to. With nothing holding
them, you compose both sides yourself, which is the whole of the difference
and is two commands rather than one.

\index{schema check!the half that stays hard}%
What that does not give you is the second comparison, and I want to be exact
about how far short it falls. The diff tells you which of the sixteen
categories you are in. It cannot tell you whether it matters, and in practice
the argument about a removal is never about the category.

\index{operation corpus!is not a population of clients}%
What stands in for the second comparison, with no control plane behind you, is
a corpus of operations you keep yourself and replay through a router started on
the candidate config. Build it. It is worth having and it is not the same
thing, for the reason
section~\ref{sec:ch16-the-change-this-book-already-shipped} gave: a corpus
maintained by the people who change the schema moves when the schema moves. It
is a regression suite, and a regression suite tells you that the graph still
does what you meant. A population of clients tells you what somebody else
meant, and only one of those two is on your machine.

\index{figure!what each check can see}%
Figure~\ref{fig:ch16-what-each-check-can-see} is the three checks against the
life of one change.

\begin{figure}[htbp]
  \centering
  % One schema change, one axis, three checks measured against it: what each
% one can actually see over the change's life, and where the two that exist
% locally leave a gap that lines up exactly with the failure in the sibling
% figure ch16-two-places-a-change-can-fail.tex.
%
% Same idiom as figures/tikz/ch01-one-field.tex: x=1mm/y=1mm, a single-weight
% axis with ticks, event nodes anchored south above it, a span bracket under
% the stages a label needs to explain. One lane, not two, because there is
% one change here, not two compared side by side; the span bracket is simply
% used three times, at three depths, so the brackets that overlap in x do not
% collide in y. No cross mark appears anywhere in this figure: nothing on
% this axis is a step a walk cannot take, only three different fields of
% view over the same four stages, so ch07-the-walk.tex's mark for a blocked
% step has no work to do here.
%
% The third bracket is the one addition. Dashing already means, in
% ch07-the-walk.tex and ch15-the-order-decides.tex, a span nobody actually
% walks; this figure reuses that meaning for a check nobody has written.
% It is also left open at the right, with no closing tick, rather than drawn
% as a complete bracket: a closed bracket over stages 1 to 2 would say a
% check exists that covers exactly that span, and it does not. The check
% this label describes would need to see stage 4, the operations clients
% actually send, while it is standing at stage 1, and nothing between those
% two points is on this axis today. An open end says the span is unclosed,
% which is the whole claim; closing it, even with a different mark, would
% say the opposite.
%
% No quotation marks anywhere in this file, including this comment block,
% per the note at the head of ch07-the-walk.tex and repeated in
% ch15-the-order-decides.tex.
%
% Bare tikzpicture (house style): the figure environment, caption and label
% live at the call site in the section file.
\begin{tikzpicture}[
  x=1mm, y=1mm,
  every node/.append style={font=\sffamily\scriptsize},
  axis/.style={draw, line width=0.4pt,
    -{Straight Barb[length=1.4mm, width=1.6mm]}},
  tick/.style={draw, line width=0.4pt},
  event/.style={anchor=south, align=center, text width=28mm, inner sep=1pt},
  span/.style={draw, line width=0.4pt},
  spanlabel/.style={anchor=north, align=center, inner sep=2pt},
  gap/.style={draw, line width=0.4pt, dashed},
]

\draw[axis] (0,0) -- (146,0);
\foreach \x in {15,54,93,131}{\draw[tick] (\x,-1.5) -- (\x,1.5);}

\node[event] at (15,3)  {you edit your\\subgraph document};
\node[event] at (54,3)  {composition merges\\four documents};
\node[event] at (93,3)  {the router loads\\the config};
\node[event, text width=32mm] at (131,3)
  {a client sends an\\operation it wrote\\last year};

% ------------------------------------------------- the two checks that exist
\draw[span] (47,-4) -- (47,-7) -- (61,-7) -- (61,-4);
\node[spanlabel, text width=40mm] at (54,-8)
  {composition. Reads the four documents and nothing else, so every
   question it can ask is a question about them};

\draw[span] (124,-4) -- (124,-7) -- (138,-7) -- (138,-4);
\node[spanlabel, text width=40mm] at (131,-8)
  {the router validates the operation. Correct, and after the
   deployment, in somebody else's request};

% ------------------------------------------- the check nobody has written
% Dashed throughout, and open at the right: no closing tick at stage 2,
% because the span this check would need does not close on this axis.
\draw[gap] (15,-22) -- (15,-25);
\draw[gap] (15,-25) -- (54,-25);
\node[spanlabel, text width=48mm] at (34,-26)
  {the check that would have caught it. Needs the operation on the right of
   this axis while standing at the edit on the left, and nothing in this book
   has it};

\end{tikzpicture}

  \caption{Three checks over the life of one schema change, and the two that
  run locally cover the wrong parts of it. Composition reads the four
  documents when they merge, and can ask nothing outside them. The router
  validates the operation correctly and does it at the far end, after the
  deployment, in somebody else's request. The third span is drawn open at its
  right because nothing in this book closes it: a check that stopped the edit
  would have to know, while you are still editing, what a client is going to
  send. A control plane is a service that closes that span by keeping the
  operations for you.}
  \label{fig:ch16-what-each-check-can-see}
\end{figure}

\subsection{What I Would Do}

\index{breaking change!my judgment on the risk}%
I think a breaking change is the largest availability risk a federated graph
carries, larger than a subgraph outage and larger than anything in
chapter~\ref{ch:blame-routing}. I state that as my judgment, because I looked
for a published account of a federated graph taken down by a schema change and
found none from anybody running one at scale. The reasoning is the one this
chapter measured. An outage ends when the service comes back, and
chapter~\ref{ch:null-blast-radius} showed the graph degrading around one rather
than emptying. A breaking change does not end. It is deployed, it composes,
it serves, and every client that meets it is broken until somebody ships a
release, which for a mobile application is a number of weeks decided by
somebody else.

\index{breaking change!the vendor who agrees, and what he sells}%
Stefan Avram, a co-founder of WunderGraph, calls it \enquote{the scariest
failure mode}, a field removed or a type renamed or an argument made
required, slipping into production and taking a client
down~\autocite{avram2026federation}. He is right and he sells the product that
catches it, which is why his sentence is here as agreement rather than as
evidence.

\index{schema check!what to run, in order}%
So, concretely, in the order I would add them. Diff the composed schema on
every change, because it costs two compositions and it turns a class of
mistake into a line of output. Keep an operation corpus and run it through a
router on the candidate config, knowing exactly what it is: your operations,
not your clients'. And when the graph has enough clients that you cannot name
them, that is the point at which the thing a control plane sells you, a record
of what was actually asked for, stops being a convenience. It is the only way
anyone has to answer the question, and everything in this chapter is evidence
that the question does not answer itself.
\index{schema check!the tool that exists|)}%
